\begin{textsample}{POS Dim 2 – global\_north\_2024\_06 – Score 19.00 – global\_north\_2024\_06/109445964.txt}  \label{ex:f2_pos_319}
Away from home, Israeli evacuees wait as Hezbollah \textbf{tensions} spike The Gil family lived so close to the \textbf{Lebanese} \textbf{border} that the sound of \textbf{rockets} often arrived before sirens ( SHARON ARONOWICZ ) Yarden Gil opens a reinforced metal door to enter the northern Israeli kindergarten where she works, which doubles as an underground shelter against \textbf{rockets} fired by Lebanon ' s Hezbollah movement.
She is among tens of thousands displaced from the \textbf{border} area by the ever-present threat of Hezbollah \textbf{attacks} and, increasingly, the fear of an \textbf{all-out} war against the powerful \textbf{Iran-backed} \textbf{militant} group.
Gil, 36, and her family have left their home in Yiftah, a kibbutz community just a few hundred metres ( yards ) from the \textbf{Lebanese} \textbf{border}.
She said there they lived so close to the \textbf{border} that they could often hear \textbf{incoming} \textbf{rockets} before the sirens started wailing.
They now live in a single room in a hotel 50 kilometres ( 30 miles ) to the south, near the city of Tiberias on the shores of the lake known as the Sea of Galilee. @ @ @ @ @ @ @ @ @ @ said Gil, charging that the Israeli government is"not doing enough for us to be able to go back to our home and be secure".
Dozens of northern Israeli communities have been rendered ghost towns as the Israeli military and Hezbollah have traded \textbf{near-daily} \textbf{cross-border} fire, ending a period of relative calm since a 2006 war.
The spike in violence during the ongoing Gaza conflict has re-ignited fears of a \textbf{wider} war between long-term \textbf{foes} Israel and Hezbollah, a Hamas \textbf{ally}.
The \textbf{border} \textbf{clashes} have killed at least 93 civilians in Lebanon and nearly 390 others, mostly \textbf{fighters}, according to an AFP tally.
Eleven civilians and 15 soldiers have been killed on the Israeli side, according to the military. - Constant ' instability ' - Israel said early last week it had approved military plans for an offensive in southern Lebanon.
Hezbollah ' s leader Hassan Nasrallah responded with a warning that nowhere in Israel would be safe in the event of war.
With Israel focused on the Gaza @ @ @ @ @ @ @ @ @ @ return home is all that is on the minds of evacuees from northern communities languishing in hotels turned state-funded shelters, away from home.
The authorities have repeatedly extended accommodation arrangements, which are now set to expire in August.
Some evacuees have moved out of the hotels, to elsewhere in Israel or abroad."That ' s our new reality : instability,"said Iris Amsalem, a 33-year-old mother of two from the \textbf{border} community of Shomera who is now staying in a Galilee hotel."We want peace.
We want security."Only a few Israelis have remained on the \textbf{border}, defended by civilian units and military forces.
Deborah Fredericks, an 80-year-old retiree staying at a five-star hotel with hundreds of other evacuees, played the tile-based game of Rummikub next to a gleaming pool and palm trees in front of the lake."It ' s really funny because I ' m in the middle of a war but I ' m on holiday,"she said."@ @ @ @ @ @ @ @ @ @ be for a while.
It ' ll be when they say I can.
You can't do anything about it."- ' War must happen ' - Others feel they have been abandoned by Prime Minister Benjamin Netanyahu ' s government as it prioritises the Gaza war."No one communicates with us, no one!
No one came to see us!"said Lili Dahn, a resident of the \textbf{border} town of Kiryat Shmona, in her 60s.
Gil, the kindergarten teacher, said parents had to set up their own schooling for their children after they fled their kibbutz, which has suffered damage from \textbf{rockets} and in fires caused by the \textbf{strikes}."The government is responsible for our security and I expect them to be more interested in what happened to us,"she said, adding that some of her fellow kibbutzniks have moved as far away as Canada and Thailand.
Netanyahu has pledged to return security, and civilians, to the north.
Some evacuees @ @ @ @ @ @ @ @ @ @ matter of time.
Sarit Zehavi, a former Israeli army intelligence official who lives near the \textbf{border}, said her greatest fear was that a potential ceasefire would allow Hezbollah"to preserve its \textbf{capabilities} and \textbf{launch} the next massacre", like Hamas did.
Gil ' s husband, Edward, 39, also said he feared a similar assault to the October 7 \textbf{attack} on southern Israel."It happened in the south,"he said."Who ' s telling me that now it won't happen in the north?"Helene Abergel, a 49-year-old Kiryat Shmona resident who is living at a Tel Aviv hotel, said :"A war must happen to push Hezbollah away from the \textbf{border}."

% matched lemmas: all-out, ally, attack, border, capability, clash, cross-border, fighter, foe, incoming, iran-backed, launch, lebanese, militant, near-daily, rocket, strike, tension, wide
\end{textsample}
